% !TEX root = chapter-standalone.tex

\section{Stacking categories}
\label{sec:parallel-stacking}

\todotext{This is strange here, move to intro?}
So far we have seen how we can compose morphisms ``horizontally'':
%
\begin{equation}
    \prfperiod{
        \mora \colon \Obja \mto \Objb
    }{
        \morb\colon \Objb \mto \Objc
    }{
        \morab \colon \Obja \mto \Objc
    }
\end{equation}
%
There are other notions of composition that allow us to compose morphisms by ``stacking them vertically''.
Given two morphisms
%
\begin{align}
    \mora & \colon \Obja \mto \Objb, \\
    \morb & \colon \Objc \mto \Objd,
\end{align}
%
we will obtain by parallel stacking a morphism
%
\begin{align}
    (\mora\mtimescatmor \morb)
     & \colon (\Obja \mtimescatob \Objc) \mto (\Objb \mtimescatob \Objd),
\end{align}
%
where ``$\mtimescatmor$'' and ``$\mtimescatob$'' are operations to be defined.
Note that while in the case of sequential composition there was a compatibility condition to be defined, as the target of the first morphism must be the source of the second morphism, here instead we can stack arbitrary morphisms.

We also expect to be able to stack any number of morphisms.
Having a collection of morphisms
%
\begin{equation}
    \moran{i} \colon \Objan{i} \mto \Objbn{i}, \qquad 1 \leq i \leq n,
\end{equation}
%
we expect to be able to obtain the composed morphism
%
\begin{equation}
    (\mtimescatmor_{i=1}^n \moran{i} )
    \colon( \mtimescatob_{i=1}^n \Objan{i} )\mto (\mtimescatob_{i=1}^n \Objbn{i}).
\end{equation}

\begin{marginfigure}
    \centering
    \includesag{parallelism-stacked-morphisms}
    %\includegraphics[scale=0.12]{parallelism_stacked-morphisms}
    \caption{Stacked morphisms}
    \label{fig:stacked-morphisms}
\end{marginfigure}

\begin{marginfigure}
    \centering
    \includesag{parallelism_stacked-morphisms-diagram}
    %\includegraphics[scale=0.17]{parallelism_stacked-morphisms-diagram}
    \caption{Stacking string diagrams}
    \label{fig:stacked-morphisms-diagram}
\end{marginfigure}

\begin{ctdefinition}[Stacking category]
    \label{def:simple-stacking-cat}
    A \maindef{stacking category} is a \SY{category}~\CatC with the following additional constituents and properties.

    \constit
    \begin{itemize}
        \item A stacking operation~$\mtimescatob \colon \ObC \cartprod \ObC \sto \ObC$.
        \item A stacking operation~$\mtimescatmor \colon \Mor_\CatC \cartprod \Mor_\CatC \sto \Mor_\CatC$.
    \end{itemize}

    \condit
    \begin{itemize}
        \item The two operations~$\mtimescatob$ and~$\mtimescatmor$ are compatible in the sense that
              %
              \begin{equation}
                  \prfperiod{
                      \moran{1} \colon \Objan{1} \mto \Objbn{1}
                      \qquad
                      \moran{2} \colon \Objan{2} \mto \Objbn{2}
                  }{
                      \moran{1} \mtimescatmor \moran{2} \colon  \Objan{1} \mtimescatob  \Objan{2} \mto \Objbn{1} \mtimescatob \Objbn{2}
                  }
              \end{equation}
    \end{itemize}
\end{ctdefinition}

In \cref{fig:stacked-morphisms} we have depicted a string diagram of two stacked morphisms.
Alternatively, in \cref{fig:stacked-morphisms-diagram} we depict the stacking of the string diagrams for $\moran{1}$ and $\moran{2}$, respectively, by stacking their diagrams vertically and drawing a box around them, merging their respective input and output terminals.
The outer box denotes $\moran{1} \mtimescatmor \moran{2}$; we think of \cref{fig:stacked-morphisms} as a ``black-boxed'' version of \cref{fig:stacked-morphisms-diagram}.

\begin{figure*}[h!]
    \centering
    \begin{subfigure}{0.49\linewidth}
        \includesag{parallelism_associative-stacking-1}
        %\includegraphics[scale=0.15]{parallelism_associative-stacking-1}
        \caption{One way of stacking three morphisms}
        \label{fig:associative-stacking-1}
    \end{subfigure}
    %
    \begin{subfigure}{0.49\linewidth}
        \includesag{parallelism_associative-stacking-2}
        %\includegraphics[scale=0.15]{parallelism_associative-stacking-2}
        \caption{Another way of stacking three morphisms}
        \label{fig:associative-stacking-2}
    \end{subfigure}
    \caption{Stacking three morphisms.}
\end{figure*}

\begin{example}
    The cartesian product of sets and functions
    \begin{equation}
        \setA \mtimescatob \setB = \setA \cartprod \setB
        \qquad \qquad
        \mapa \mtimescatmor \mapb = \mapa \funcprod \mapb
    \end{equation}defines a stacking operation on the category $\Set$.
    Indeed, by the definition of the cartesian product of functions, the two stacking layers are compatible:
    \begin{equation}
        \prfperiod{
            \mapa \colon \setA \to \setC
        }{\qquad}{
            \mapb \colon \setB  \to \setD
        }{
            \defmap{
                \mapa \funcprod \mapb
            }{
                \setA \cartprod \setB
            }{
                \mto
            }{
                \setC \cartprod \setD
            }{
                \tup{\setAel, \setCel}
            }{
                \tup{\mapa(\setAel),\mapb(\setCel)}
            }
        }
    \end{equation}
\end{example}

\begin{example}
    The sum of sets and functions
    \begin{equation}
        \setA \mtimescatob \setB = \setA \setdisunion \setB
        \qquad \qquad
        \mapa \mtimescatmor \mapb = \mapa \funcsum \mapb
    \end{equation}
    also defines a stacking operation on the category $\Set$.
\end{example}

\todotext{J: product and sum for Rel -- make exercise or graded exercise perhaps?}

\begin{example}
    The category of real matrices admits a stacking operation defined by summing dimensions
    \begin{equation}
        n \mtimescatob m = n + m
    \end{equation}
    and combining block matrices like this:
    \begin{equation}
        \mat{A} \mtimescatmor \mat{B} = \begin{bmatrix} \mat{A} & \mat{0} \\ \mat{0} & \mat{B} \end{bmatrix}.
    \end{equation}
\end{example}

\begin{example}
    The category $\DP$ of design problems admits a stacking operation which is defined on objects by taking the product of posets,
    \begin{equation}
        \posA \mtimescatob \posB = \posA \Ptimes \posB,
    \end{equation}
    and on morphisms it is defined by
    \begin{equation}
        \prfperiod{
            \adpa\colon \posAop\Ptimes \posC \mto \Bool
        }{\qquad}{
            \adpb\colon \posB\posop\Ptimes \posD \mto \Bool,
        }{
            \defmap{
                \adpa \mtimescatmor \adpb
            }{
                \pars{\posA \mtimescatob \posB}\op
                \Ptimes
                \pars{\posC \mtimescatob \posD}
            }{
                \mto
            }{
                \Bool
            }{
                \tup{\tup{a, c},\ \tup{b, d}}
            }{
                \adpa(a,c) \booland \adpb(c, d)
            }
        }
    \end{equation}
\end{example}

\begin{example}
    The following defines a stacking operation for the category \LTI of LTI systems.
    On objects the stacking is just addition on the natural numbers (which represent dimensions of input and output spaces):
    \begin{equation}\label{eq:LTI-stacking-objects}
        \defmapperiod{
            \mtimescatob
        }{
            \Obof\LTI \cartprod \Obof\LTI
        }{
            \to
        }{
            \Obof\LTI
        }{
            \tup{\styleobj{l},\styleobj{m}}
        }{
            \styleobj{l}+\styleobj{m}
        }
    \end{equation}
    On morphisms, the stacking is\footnote{For the control engineers out there: the resulting LTI system will have a Relative Gain Array (RGA) matrix corresponding to the identity matrix.
    }
    \begin{equation}\label{eq:LTI-stacking-prftree}
        \prfperiod{
            \mora \colon \styleobj{l}\mtoin\LTI \styleobj{m}
        }{
            \morb \colon \styleobj{n}\mtoin\LTI \styleobj{o}
        }{
            \mora\mtimescatmor\morb = \genericlti{} \colon
            \styleobj{l+n}\mtoin\LTI \styleobj{m+o}
        }
    \end{equation}
    % \begin{equation}
    %     \defmapcomma{
    %         \mtimescatmor
    %     }{
    %         \Mor_\LTI \cartprod \Mor_\LTI
    %     }{
    %         \to
    %     }{
    %         \Mor_\LTI
    %     }{
    %         \tup{\genericlti{\mora},\genericlti{\morb}}
    %     }{
    %         \genericlti{}
    %     }
    % \end{equation}
    with
    \begin{equation}\label{eq:LTI-stacking-mor1}
        \prstart=\begin{bmatrix}
            \prstart_\mora \\
            \prstart_\morb
        \end{bmatrix},\
        \mat{A}=\begin{bmatrix}
            \mat{A}_\mora & \mat{0}       \\
            \mat{0}       & \mat{A}_\morb
        \end{bmatrix},\
        \mat{B}=\begin{bmatrix}
            \mat{B}_\mora & \mat{0}       \\
            \mat{0}       & \mat{B}_\morb
        \end{bmatrix},
        %  \
        % \mat{B}=\begin{bmatrix}
        %     \mat{B}_\mora & \mat{0}       \\
        %     \mat{0}       & \mat{B}_\morb
        % \end{bmatrix},\
        % \mat{C}=\begin{bmatrix}
        %     \mat{C}_\mora & \mat{0}       \\
        %     \mat{0}       & \mat{C}_\morb
        % \end{bmatrix},\
        % \mat{D}=\begin{bmatrix}
        %     \mat{D}_\mora & \mat{0}       \\
        %     \mat{0}       & \mat{D}_\morb
        % \end{bmatrix}
    \end{equation}
    \begin{equation}\label{eq:LTI-stacking-mor2}
        % \prstart=\begin{bmatrix}
        %     \prstart_\mora \\
        %     \prstart_\morb
        % \end{bmatrix},\
        % \mat{A}=\begin{bmatrix}
        %     \mat{A}_\mora & \mat{0}       \\
        %     \mat{0}       & \mat{A}_\morb
        % \end{bmatrix}, \
        % \mat{B}=\begin{bmatrix}
        %     \mat{B}_\mora & \mat{0}       \\
        %     \mat{0}       & \mat{B}_\morb
        % \end{bmatrix},\
        \mat{C}=\begin{bmatrix}
            \mat{C}_\mora & \mat{0}       \\
            \mat{0}       & \mat{C}_\morb
        \end{bmatrix},\
        \mat{D}=\begin{bmatrix}
            \mat{D}_\mora & \mat{0}       \\
            \mat{0}       & \mat{D}_\morb
        \end{bmatrix}.
    \end{equation}

    % Clearly, the morphism given by~$\genericlti{}$ satisfies the compatibility condition for stacking semicategories:
    % if~$\genericlti{\mora}$ is~$\styleobj{l}\mto \styleobj{m}$ and~$\genericlti{\morb}$ is~$\styleobj{n}\mto \styleobj{o}$,~$\genericlti{}$ is~$\styleobj{l}+\styleobj{n}\mto \styleobj{m}+\styleobj{o}$.
\end{example}

\todotext{Put examples here of stacking categories that are in the slides}

